% Generated draft: requires mathematical and source review.
\documentclass[11pt]{article}
\usepackage[margin=25mm]{geometry}
\usepackage{fontspec}
\setmainfont{Latin Modern Roman}
\usepackage{amsmath,amssymb,mathtools,unicode-math}
\setmathfont{Latin Modern Math}
\usepackage{enumitem,needspace,xurl,hyperref}
\setlength{\parindent}{0pt}
\setlength{\parskip}{5pt}
\setlength{\emergencystretch}{3em}
\newcommand{\N}{\mathbb{N}}
\newcommand{\Z}{\mathbb{Z}}
\newcommand{\Q}{\mathbb{Q}}
\newcommand{\R}{\mathbb{R}}
\newcommand{\C}{\mathbb{C}}
\newcommand{\F}{\mathbb{F}}
\newcommand{\A}{\mathbb{A}}
\newcommand{\PP}{\mathbb{P}}
\newcommand{\E}{\mathbb{E}}
\newcommand{\eps}{\varepsilon}
\newcommand{\dd}{\,\mathrm{d}}
\begin{document}
\section*{UnsolvedMath ID: 30005312}
\textbf{Source problem code:} catalogue entry
\subsection{Definitions and mathematical statement}
A \emph{scriptor} here means a person who writes the notes or manuscript under study. The supplied passage distinguishes at least two voices in reading or lecture notes: that of the note-taker and that of the text or speaker being recorded. It attributes the expression ``several enunciative instances'' to its reference [18], whose bibliographic details are not supplied. A \emph{private manuscript}, in the sense used in the passage, is a manuscript intended only for its author's eyes. The passage observes that such manuscripts often depart from customary writing norms and suggests that methods from genetic criticism may help their analysis.

\textbf{Formulation limitation.} The supplied material presents methodological questions in the history of mathematics, rather than a mathematical conjecture or a theorem to be proved. It specifies neither a particular manuscript corpus nor mathematical objects, hypotheses, quantifiers, or a conclusion. Consequently, no precise mathematical statement can be recovered without adding information absent from the source. The following preserves the historical research questions instead of inventing a mathematical formulation.

For the reading notes, lecture notes, notebooks, slips of paper, and loose sheets considered in the passage, investigate the following questions using material evidence together with analysis of their content:
\begin{enumerate}
\item Who is writing? Are several scriptors involved? Are there changes in pens or handwriting? How can the note-taker's voice be distinguished from that of the text or speaker being recorded?
\item What is the lecturer's status, and what are the scriptor's objectives?
\item Which writing acts can be observed, including underlining, titles, numbering, textual modifications, rewriting, and annotations?
\item Are notebooks organized thematically, chronologically, or according to another discernible arrangement?
\item When was a manuscript written? Reused paper and watermarks may provide dating clues, but the passage explicitly cautions that these must be used carefully.
\item For whom was the manuscript written? In particular, how should the intended privacy of a manuscript inform its interpretation?
\item What is the chronology of the writing process?
\end{enumerate}

The passage situates these questions within historical studies of reading and lecture notes, while describing the history of mathematics as noticeably absent from most such studies. Its notebook examples concern Cartan; its examples involving slips and loose sheets concern Leibniz's and Dedekind's archives. It presents material analysis as a complement to content analysis, not as a guarantee of determinate answers. The supplied excerpt has a disrupted transition before its notebook discussion, so the missing wording cannot be reconstructed reliably.

The registry labels the item ``open'', but also marks its statement as \texttt{unrecoverable} and supplies the verification note that the official report passage concerns historical attribution and manuscript analysis rather than an open mathematical problem. These are supplied curation labels and qualifications, not an independently verified mathematical status assessment.
\subsection{Short English statement}
The source asks how historians can identify writers, intended readers, dates, writing operations, and stages of composition in mathematical notes and private manuscripts. It gives methodological questions rather than a mathematical conjecture, and does not identify a specific corpus to investigate. A precise mathematical statement cannot be recovered from the supplied evidence; the registry's ``open'' label must be read alongside its \texttt{unrecoverable} statement classification.
\subsection{Sources}
\begin{itemize}
\item Supplied local UnsolvedMath registry record, permanent ID 30005312, title ``Authorship and Composition of Private Mathematical Manuscripts'', source problem code \texttt{catalogue entry}. Project registry locations: \texttt{data/unsolvedmath/problems.json} and the supplemental \texttt{catalog/problems.json}. This draft uses the supplied record without independently inspecting those files. No proposer is identified; the record gives proposed year 2022.
\item Supplied report citation: \emph{History of mathematics through collaboration: Toward a composite portrait of Oswald Veblen} (2023), \emph{Oberwolfach Reports}, pp.\ 3171--3214, DOI 10.4171/owr/2022/56. \url{https://doi.org/10.4171/owr/2022/56}. The same link is supplied under the literature-source label ``original OWR report''. The 2023 publication attribution and the record's 2022 year designation are retained as supplied.
\item Supplied literature source: ``Library of Congress primary finding aid''. \url{https://findingaids.loc.gov/repositories/19/resources/4615}. The supplied record does not identify the collection by name or establish its connection to particular manuscripts in the excerpt.
\item Internal references retained from the supplied excerpt: [18] for ``several enunciative instances''; [12] for Cartan; [1, 16, 15] for additional notebook-related insights; [13, 14] for Dedekind; [10] for Leibniz's mathematical manuscripts; [11] for genetic criticism; [17], attributed to Leong, for an example concerning writing chronology; and [2, 4, 5, 6, 7, 18, 19] for studies of reading and lecture notes. Their bibliographic entries are not supplied, so these labels cannot be resolved here. The excerpt associates footnotes 3--6 with these discussions.
\item Supplied literature assessment, checked at \texttt{2026-08-22}: status ``open''; the record characterizes the questions as methodological and states that archival sources describe the collection without resolving the listed authorship, chronology, or writing-operation questions. This is supplied dated curation derived from the OWR record, not an independently verified assessment. The record itself calls for independent verification of time-sensitive status and cited claims.
\end{itemize}
\end{document}
